%============================================================
% how_to_use.tex — reading workflow, chapter anatomy, legend
%============================================================
\chapter{How to Use This Book}

\section*{The Anatomy of Every Chapter}

Every chapter has the same skeleton, so you always know where you are.

\begin{diagrambox}
\begin{tikzpicture}[node distance=4mm and 5mm,
    blk/.style={rectangle, rounded corners=3pt, draw=#1, fill=#1!10, line width=1pt,
                align=center, font=\small, minimum height=10mm, text width=2.5cm},
    prb/.style={rectangle, rounded corners=2pt, draw=primaryblue, fill=white, line width=0.8pt,
                align=center, font=\scriptsize, minimum height=6mm, text width=1.05cm}]
    \node[blk=theoryborder] (T) {\textbf{Part A}\\{\scriptsize 5 lessons}};
    \node[blk=patternborder, right=of T] (P1) {\textbf{Pattern 1}\\{\scriptsize theory}};
    \node[blk=patternborder, right=of P1] (P2) {\textbf{Pattern 2}\\{\scriptsize theory}};
    \node[right=2mm of P2, font=\Large\bfseries\color{patternborder}] (dots) {$\cdots$};
    \node[blk=patternborder, right=2mm of dots] (P5) {\textbf{Pattern 5}\\{\scriptsize theory}};
    \foreach \p in {P1,P2,P5} {
        \node[prb, below=5mm of \p, xshift=-6.5mm] (\p a) {problem 1};
        \node[prb, below=5mm of \p, xshift=6.5mm] (\p b) {problem 2};
        \node[prb, below=2.5mm of \p a, xshift=6.5mm, text width=2.35cm, draw=successgreen] (\p e) {exercises};
        \draw[faded] (\p) -- (\p a); \draw[faded] (\p) -- (\p b);
    }
    \node[blk=successgreen] (S) at (T |- P1e) {\textbf{Summary}\\{\scriptsize + practice set (30+)}};
    \draw[flow] (T) -- (P1); \draw[flow] (P1) -- (P2); \draw[flow] (P2) -- (dots); \draw[flow] (dots) -- (P5);
    \draw[flow] (P5e.south) -- ++(0,-0.45) -| (S.south);
\end{tikzpicture}
\figcap{Part A lessons $\to$ 5 patterns (each: theory, 2 fully solved problems, exercises) $\to$ summary $\to$ practice set with solution sketches.}
\end{diagrambox}

\section*{Workflow for Each Solved Problem}

\begin{insightbox}[The 5-Step Ladder --- climb it for every solved problem]
\begin{center}
\begin{tikzpicture}[node distance=3.5mm,
    stepbox/.style={rectangle, rounded corners=3pt, line width=1pt, align=center,
                    font=\small\bfseries, minimum height=11mm, text width=2.15cm}]
    \node[stepbox, draw=definitionborder, fill=definitionbg] (m) {1. Model\\{\scriptsize\mdseries vertices, edges, weights}};
    \node[stepbox, draw=recursionred, fill=recursionbg, right=of m] (b) {2. Brute force\\{\scriptsize\mdseries correct but slow}};
    \node[stepbox, draw=tabulationblue, fill=tabulationbg, right=of b] (a) {3. Algorithm\\{\scriptsize\mdseries the right tool}};
    \node[stepbox, draw=deepblue, fill=white, right=of a] (p) {4. Prove\\{\scriptsize\mdseries the invariant}};
    \node[stepbox, draw=secondaryblue, fill=white, right=of p] (t) {5. Trace\\{\scriptsize\mdseries run it by hand}};
    \draw[flow] (m) -- (b); \draw[flow] (b) -- (a); \draw[flow] (a) -- (p); \draw[flow] (p) -- (t);
\end{tikzpicture}
\end{center}
\begin{enumerate}[leftmargin=*]
    \item Read the \textbf{Part A lessons} once quickly, then again slowly after the patterns.
    \item For each \textbf{pattern}: study the framework, then \emph{try} the two solved problems yourself for 20--30 minutes before reading the solutions.
    \item Start every problem with the \textbf{Graph Model}: what is a vertex, what is an edge, is it directed, weighted, implicit? Most wrong answers in graph problems are wrong models, not wrong algorithms.
    \item Redraw every trace on paper; then read the \textbf{Variants} box to see which other problems the same model solves.
    \item Finish the chapter's \textbf{practice set}; do not move on until the checklist in the summary is ticked.
\end{enumerate}
\end{insightbox}

\begin{insightbox}[Companion book]
Graph algorithms and dynamic programming meet often: DP on DAGs, DP on trees, rerooting and bitmask DP over vertex sets. Those techniques live in the companion book \textit{Dynamic Programming --- From Foundations to World Finals}; this book points to it (``\dpbook{11}{DP on Trees}'') instead of repeating it, and no problem is solved in both books.
\end{insightbox}

\clearpage
\section*{Legend}

Three kinds of colour carry meaning in this book: the \textbf{boxes} around text, the \textbf{badges} next to problem names, and the \textbf{vertices and edges} of graph diagrams.

%------------------------------------------------------------
\begin{tcolorbox}[enhanced, colback=white, colframe=deepblue, boxrule=0.8pt, arc=3pt,
    title={\bfseries 1\quad Boxes}, colbacktitle=deepblue, fonttitle=\small, left=6pt, right=6pt]
\centering
\begin{tikzpicture}[sw/.style 2 args={rectangle, rounded corners=2pt, draw=#1, fill=#2, line width=1.2pt,
                    minimum width=9mm, minimum height=5.5mm},
                    lab/.style={anchor=west, font=\small, text width=3.7cm}]
    \foreach \bd/\bg/\txt [count=\i from 0] in {
        theoryborder/theorybg/Chapter theory, patternborder/patternbg/Pattern theory,
        definitionborder/definitionbg/Graph model (and definitions),
        recursionred/recursionbg/Approach 1: brute force, tabulationblue/tabulationbg/Approach 2: efficient algorithm,
        memoorange/memobg/Approach 3: alternative, optimizedpurple/optimizedbg/Approach 4: refinements,
        softpurple/softpurple!4/Variants of a solved problem} {
        \node[sw={\bd}{\bg}] at (0,-\i*0.72) {};
        \node[lab, text width=5.4cm] at (0.6,-\i*0.72) {\txt};
    }
    \foreach \bd/\bg/\txt [count=\i from 0] in {
        insightborder/insightbg/Key insight, warningborder/warningbg/Common mistake,
        deepblue/white/Invariant or theorem, secondaryblue/white/Algorithm trace,
        solutiongreen/solutionbg/Solution and tests, primaryblue/white/Practice problem,
        gray/lightgray/Diagram, midgray/lightgray/Planned chapter} {
        \node[sw={\bd}{\bg}] at (7.0,-\i*0.72) {};
        \node[lab, text width=5.0cm] at (7.6,-\i*0.72) {\txt};
    }
\end{tikzpicture}
\end{tcolorbox}

%------------------------------------------------------------
\begin{tcolorbox}[enhanced, colback=white, colframe=deepblue, boxrule=0.8pt, arc=3pt,
    title={\bfseries 2\quad Difficulty badges}, colbacktitle=deepblue, fonttitle=\small, left=6pt, right=6pt]
\renewcommand{\arraystretch}{1.7}
\begin{tabularx}{\linewidth}{@{}>{\raggedright\arraybackslash}p{3.0cm} X@{}}
\textbf{Book badge}\newline{\scriptsize\color{darkgray}where the problem sits in its chapter} &
\mbox{\exeasy\ {\small warm-up}}\hspace{4mm plus 2mm} \mbox{\exmedium\ {\small core}}\hspace{4mm plus 2mm} \mbox{\exhard\ {\small challenge}}\hspace{4mm plus 2mm} \mbox{\exstretch\ {\small above the chapter}}\hspace{4mm plus 2mm} \mbox{\exexpert\ {\small expert}} \\
\textbf{Codeforces rating}\newline{\scriptsize\color{darkgray}in Codeforces rank colours} &
\cf{1000}\ \cf{1300}\ \cf{1500}\ \cf{1700}\ \cf{2000}\ \cf{2300}\ \cf{2600} \\
\textbf{LeetCode label}\newline{\scriptsize\color{darkgray}LeetCode's own difficulty} &
\lc{Easy}\ \ \lc{Medium}\ \ \lc{Hard} \\
\end{tabularx}
\par\smallskip
{\small\color{darkgray}The book badge and the judge's label can differ: Word Ladder is \lc{Hard} on LeetCode, but only \exmedium{} once you know multi-source and bidirectional BFS. Every problem also has a \textbf{Solve} button that opens it on the judge.}
\end{tcolorbox}

%------------------------------------------------------------
\begin{tcolorbox}[enhanced, colback=white, colframe=deepblue, boxrule=0.8pt, arc=3pt,
    title={\bfseries 3\quad Vertices and edges in diagrams}, colbacktitle=deepblue, fonttitle=\small, left=6pt, right=6pt]
\centering
\begin{tikzpicture}[lab/.style={anchor=west, font=\small}]
    \foreach \st/\txt [count=\i from 0] in {gvsrc/source, gvcur/being processed, gvseen/discovered (queued)} {
        \node[\st] at (\i*4.4,0) {};
        \node[lab] at (\i*4.4+0.5,0) {\txt};
    }
    \foreach \st/\txt [count=\i from 0] in {gvdone/finished, gvdst/target or answer, gvdead/removed or unreachable} {
        \node[\st] at (\i*4.4,-0.9) {};
        \node[lab] at (\i*4.4+0.5,-0.9) {\txt};
    }
    \foreach \st/\txt [count=\i from 0] in {gvA/class A (left side), gvB/class B (right side), gvsel/selected (matching or cover)} {
        \node[\st] at (\i*4.4,-1.8) {};
        \node[lab] at (\i*4.4+0.5,-1.8) {\txt};
    }
    \foreach \st/\txt [count=\i from 0] in {gtree/tree edge, gback/back edge, gfwd/forward edge} {
        \draw[\st] (\i*4.4-0.4,-2.8) -- (\i*4.4+0.4,-2.8);
        \node[lab] at (\i*4.4+0.55,-2.8) {\txt};
    }
    \foreach \st/\txt [count=\i from 0] in {gcross/cross edge, gpath/chosen path, gcut/cut} {
        \draw[\st] (\i*4.4-0.4,-3.6) -- (\i*4.4+0.4,-3.6);
        \node[lab] at (\i*4.4+0.55,-3.6) {\txt};
    }
\end{tikzpicture}
\end{tcolorbox}
